% \section{Wybrane zagadnienia teoretyczne z zakresu automatyzacji planownia zadań i personalizacji}\label{section_theory}
% \subsection{Podstawowe pojęcia automatyzacji planowania zadań}\label{section_theory_dictionary_tasks}
% - pareto front

% - human-in-the-loop

% - Problem harmonogramowania zadań (Job-Shop Scheduling / Flexible Job-Shop Problem)

% - Modelowanie ograniczeń

% \subsection{Niezdominowany algorytm sortowania genetycznego}\label{section_theory_NSGA}
% - Optymalizacja wielokryterialna

% - NSGA-II

% - NSGA-III

% \subsection{Uczenie przez wzmacnianie}\label{section_theory_RL}
% - Markov Decision Process

% - Dynamiczne przeplanowanie

% - PPO - proximal policy optimization

% - Uczenie ze wzmocnieniem z informacji zwrotnych od ludzi

% \subsection{Rytm dobowy i jego wpływ na energię i produktywność człowieka}\label{section_theory_dictionary_tasks}
\section{Wybrane zagadnienia teoretyczne z zakresu automatyzacji
planowania zadań i personalizacji}

\subsection{Problem automatycznego planowania zadań}

    \subsubsection{Planowanie i harmonogramowanie zadań}

    \subsubsection{Problemy Job-Shop Scheduling i Flexible Job-Shop Scheduling}

    \subsubsection{Modelowanie zadań, zasobów i ograniczeń}

    \subsubsection{Dynamiczne przeplanowanie zadań}

    \subsubsection{Personalizacja i podejście Human-in-the-Loop}


\subsection{Optymalizacja wielokryterialna}

    \subsubsection{Problem optymalizacji wielokryterialnej}

    \subsubsection{Dominacja Pareto i front Pareto}

    \subsubsection{Algorytmy ewolucyjne w optymalizacji wielokryterialnej}

    \subsubsection{Algorytm NSGA-II}

    \subsubsection{Algorytm NSGA-III}


\subsection{Uczenie przez wzmacnianie}

    \subsubsection{Podstawy uczenia przez wzmacnianie}

    \subsubsection{Proces decyzyjny Markowa}

    \subsubsection{Funkcja nagrody i polityka agenta}

    \subsubsection{Proximal Policy Optimization}

    \subsubsection{Informacja zwrotna użytkownika w procesie uczenia}


\subsection{Modelowanie zmienności funkcjonowania człowieka}

    \subsubsection{Rytm okołodobowy}

    \subsubsection{Chronotyp i różnice indywidualne}

    \subsubsection{Zmęczenie i wydajność podczas wykonywania zadań}

    \subsubsection{Odpoczynek i regeneracja}

    \subsubsection{Reprezentacja stanu użytkownika w systemie planowania}